% --- Sprache und Kodierung ---
\usepackage[utf8]{inputenc}
\usepackage[T1]{fontenc}
\usepackage[ngerman]{babel}
\makeatletter
\newif\ifmodernlatex
\@ifl@t@r\fmtversion{2025-06-01}{\modernlatextrue}{\modernlatexfalse}
\makeatother

% --- Layout ---
\usepackage[a4paper,margin=2.5cm]{geometry}
\usepackage{setspace} % Zeilenabstand
\onehalfspacing
\usepackage{placeins} % Floatbarrier (Kapitel)

% --- Schriftarten ---
\usepackage{lmodern}

% --- Mathe und Symbole ---
\usepackage{amsmath, amssymb, amsthm}
\usepackage{calc}

% --- Grafiken ---
\usepackage{graphicx}
\ifmodernlatex
    \usepackage{tcolorbox}
    \tcbset{
        myboxstyle/.style={
            colframe=blue!60!black,
            colback=blue!5,
            boxrule=1pt,
            arc=2mm,
            left=3pt,right=3pt,top=3pt,bottom=3pt,
            width=\textwidth,
            fontupper=\footnotesize
        }
    }
\else
    \newenvironment{tcolorbox}[1][]
        {\par\medskip\begin{quote}\small}
        {\end{quote}\par\medskip}
\fi
\usepackage{float} % bessere Kontrolle von Abbildungen
\graphicspath{{figures/}}


% --- Tabellen ---
\usepackage{longtable}
\usepackage{booktabs}
\usepackage{threeparttable}
\usepackage{tabularx}
\usepackage{makecell}
\usepackage{etoolbox}                
\AtBeginEnvironment{tabular}{\footnotesize}
\AtBeginEnvironment{tabularx}{\footnotesize}
\ifmodernlatex
    \usepackage{siunitx}
    \sisetup{
        detect-all,
        table-format=3.2,
        separate-uncertainty = true,
        output-decimal-marker = {.},
        group-separator = {'}
    }
    \DeclareSIUnit{\gon}{gon}
\else
    \newcommand{\SI}[2]{\ensuremath{#1\,\mathrm{#2}}}
    \newcommand{\gon}{gon}
    \newcommand{\meter}{m}
\fi
\usepackage{multicol}
\usepackage{multirow}
\usepackage{xcolor}


% --- Referenzen und Links ---
\usepackage{xurl}                
\usepackage{csquotes} % für saubere Zitate
\usepackage[backend=biber,style=apa]{biblatex}
\addbibresource{bibliography.bib}
\usepackage[colorlinks=true,linkcolor=black,citecolor=black,urlcolor=blue]{hyperref}
\usepackage[ngerman]{cleveref}
% Deutsche Bezeichnungen für \cref/\Cref
\crefname{figure}{Abbildung}{Abbildungen}
\Crefname{figure}{Abbildung}{Abbildungen}
\crefname{table}{Tabelle}{Tabellen}
\Crefname{table}{Tabelle}{Tabellen}
\crefname{equation}{Gleichung}{Gleichungen}
\Crefname{equation}{Gleichung}{Gleichungen}
\usepackage[labelfont=bf,font={footnotesize, it},labelsep=colon]{caption}


% --- Sonstiges ---
\setlength{\marginparwidth}{2cm}
\usepackage[textsize=tiny]{todonotes}
\usepackage{ifthen}
\usepackage{xkeyval}
\usepackage{tikz}
\usepackage{lscape}
\usepackage{pdflscape}
\usepackage{amsmath}
\usepackage{pdfpages}
\usepackage{hyperref}


\newif\iffirstpdfpage

% \anhangpdf{Kürzel}{Dateiname ohne .pdf}{Titel im Inhaltsverzeichnis}
\newcommand{\anhangpdf}[3]{%
  \clearpage
  \phantomsection
  \addcontentsline{toc}{subsection}{#3}%
  \global\firstpdfpagetrue
  \includepdf[pages=-,
    pagecommand={\iffirstpdfpage\hypertarget{pdf:#1}{}\global\firstpdfpagefalse\fi}
  ]{anhang/#2.pdf}%
}

% Verweis im Text: ergibt "Anhang 7" (anklickbar, Nummer automatisch)
\newcommand{\anhangref}[1]{\hyperlink{pdf:#1}{Anhang~\ref*{nr:#1}}}